% =============================== INTRODUCAO ==================================
% Sem subsecoes: contexto -> problema -> lacuna -> pergunta/objetivo -> contribuicoes.
\section{Introdução}\label{sec:introducao}

A cirrose biliar primária (CBP) é uma doença hepática crônica e autoimune que destrói progressivamente os pequenos ductos biliares do fígado e pode evoluir para cirrose, insuficiência hepática e necessidade de transplante. Como a maioria dos pacientes não tem cirrose no diagnóstico, a doença passou a ser chamada, a partir de 2015, de colangite biliar primária, mantendo a sigla \cite{beuers2015nomenclature}. Este artigo usa o nome antigo porque é o que consta no conjunto de dados e na literatura que o originou. A evolução da doença varia muito entre pacientes, e estimar o risco individual de óbito na primeira avaliação ajuda a definir a intensidade do acompanhamento e o momento de encaminhamento ao transplante.

Desde 1989 existe uma ferramenta para essa estimativa: o Escore de Risco de Mayo, um modelo de regressão de Cox com cinco variáveis de rotina (idade, bilirrubina, albumina, tempo de protrombina e edema), derivado de um ensaio clínico conduzido na Mayo Clinic entre 1974 e 1984 \cite{dickson1989prognosis}. Os dados desse ensaio são públicos e se tornaram um conjunto clássico de ensino e de teste de métodos, primeiro em análise de sobrevivência \cite{fleming1991counting} e, depois, em aprendizado de máquina (AM), inclusive no repositório da UCI \cite{uci878} e na competição Kaggle Playground S3E26 \cite{reade2023ps3e26}. Modelos de AM prometem mais precisão, mas costumam trazer custos: perdem a transparência de uma fórmula que o médico calcula à beira do leito, exigem mais infraestrutura e, quando usam variáveis além das cinco de Mayo, exigem também mais exames e mais pacientes com dados completos.

A pergunta relevante para a prática, portanto, não é se um modelo de AM é preciso, mas se ele é \textbf{mais} preciso que a ferramenta já disponível, e a que custo. Esse é o conceito de valor incremental: o benefício adicional que um novo modelo traz em comparação com o método clínico em uso. Uma revisão sistemática de 71 estudos não encontrou vantagem do AM sobre a regressão logística quando a comparação tinha baixo risco de viés, e observou que 79\% dos estudos não avaliaram calibração \cite{christodoulou2019systematic}. No conjunto da Mayo Clinic, o problema é agravado pelo enquadramento mais difundido da tarefa: a competição do Kaggle pede a classificação do desfecho em três classes usando o tempo de seguimento como variável de entrada, uma informação que só existe depois do desfecho e que caracteriza vazamento de dados \cite{kapoor2023leakage}.

Este trabalho pergunta quanto o AM acrescenta a um modelo de regressão com as cinco variáveis do Escore de Mayo na predição de óbito em até cinco anos, e se um eventual ganho viria do algoritmo ou das variáveis extras. O objetivo é medir esse valor incremental em discriminação e calibração, num desenho 2 x 2 que cruza a família de métodos (regressão e AM) e o conjunto de variáveis (as 5 de Mayo e as 17 do conjunto de dados), avaliado nos mesmos pacientes e nas mesmas partições.

As contribuições são quatro: (i) uma comparação controlada entre regressão e AM em que todos os modelos usam as mesmas partições, o mesmo orçamento de ajuste e transformações aprendidas só no treino; (ii) a quantificação do valor incremental com intervalos de confiança pareados e correção para comparações múltiplas, acompanhada da avaliação de calibração; (iii) a medida de quanto o enquadramento com vazamento infla o desempenho; e (iv) código e resultados abertos, com reprodução verificada, seguindo as recomendações do TRIPOD+AI \cite{collins2024tripodai}.

O restante do artigo está organizado da seguinte forma: a Seção~\ref{sec:relacionados} apresenta os trabalhos relacionados; a Seção~\ref{sec:metodos} descreve dados e procedimentos; a Seção~\ref{sec:resultados} apresenta os resultados; a Seção~\ref{sec:discussao} discute implicações e limitações; e a Seção~\ref{sec:conclusao} sintetiza as conclusões.
